\scanpage{33}
\mspage{3}

était pris comme la donnée de départ fondamentale pour faire de l'algèbre homotopique en tous genres. J'avais d'ailleurs \unsure{développé|développe} une kyrielle de notions du \unsure{$\overline{\rm m}$|$=$} genre (i) à (v), en partant de cette donnée $W$. Il se pourrait que \deleted{\unsure{toutes|tous}} les \inserted{énoncés} \deleted{\unsure{notions}} \deleted{\unsure{que|de}} sur ces notions que j'ai \unsure{développés|développé} dans $\mathrm I$, en termes d'un dérivateur, \deleted{\unsure{restent|seraient}} valables plus généralement en termes d'un localisateur $W$, \deleted{\unsure{et|je}} \unsure{dans} la mesure \unsure{du moins|des notions} où ils gardent un sens. (\unsure{Donc|Tout} pour l'instant je mets à part les énoncés proprement cohomologiques ou homologiques, qui font directement appel à $\mathbb D$. Encore qu'il se pourrait que moyennant les conditions à expliciter sur $W$, qui impliquent que $(\mathrm{Cat},W)$ définissent bel et bien un dérivateur $\mathbb D_{\mathrm{Cat},W}=\mathbb D$, on arrive encore à interpréter $W$ comme un $W_{\mathbb D}$, de sorte que les énoncés cohomologiques et homologiques gardent un sens.

J'ai envie d'abord de regarder dans PS les axiomes que j'ai dégagés pour les localisateurs, et examiner s'ils sont tous vérifiés pour
% archive sheet number: 32
